\subsection{语料边界}

本报告把名称匹配 A/B/C 加编号的 PDF 视为用户提供的国赛一等奖论文语料，不再
外推奖项名单。按年份和题型清点结果见表~\ref{tab:corpus-count}。59 篇论文全部位于
各年份目录根部，页数合计 2892。原始目录 \texttt{2020/} 至 \texttt{2025/} 及思维
导图 \texttt{1.png} 只读使用，不因 OCR、统计或排版而改写。

\begin{table}[htbp]
\centering
\caption{编号论文语料分布}
\label{tab:corpus-count}
\begin{tabular}{crrrr}
\toprule
年份 & A 题 & B 题 & C 题 & 合计\\
\midrule
2020 & 4 & 4 & 5 & 13\\
2021 & 3 & 4 & 4 & 11\\
2022 & 3 & 3 & 2 & 8\\
2023 & 4 & 3 & 3 & 10\\
2024 & 5 & 3 & 4 & 12\\
2025 & 1 & 2 & 2 & 5\\
\midrule
合计 & 20 & 19 & 20 & 59\\
\bottomrule
\end{tabular}
\end{table}

三类材料的用途不可互换。编号论文回答“获奖论文怎样组织和表达”；赛题原文回答
“当年要求解决什么”；专家评述帮助理解命题背景。后两类材料不进入获奖论文的
章节比例、段落密度、图表与公式安排或高频表述统计。2020 年本地缺少 A/B/C 赛题
原文，表~\ref{tab:problems} 只据专家评述确认题名，不据评述反推子问。

